\section{Conclusion}

Bhoonidhi solves discovery. Its STAC-based Open Data Access API is real,
modern engineering, and a developer can authenticate, search by space and
time, and download products by identifier. This report makes no complaint
about that, and none about the requirement to hold a token.

What Bhoonidhi does not do is \emph{compute}. It stores the footprint polygon
of every scene it holds and permits exactly one question of that geometry:
does it intersect mine? The audit in Section~\ref{sec:gap} establishes this
from NRSC's own published specification --- two spatial parameters expressing
one predicate, attribute equality in CQL2, and no analysis vocabulary
anywhere. The probe in Section~\ref{sec:gap} establishes that no processing
endpoint exists anywhere in ISRO's public stack: Bhuvan's WPS returns 404, as
does every modern equivalent.

Three things make that a large gap rather than a missing convenience.

\textbf{It is an absent tier.} Server-side geospatial processing has been an
OGC standard since 2007 and has a modern REST/JSON successor from 2021.
Bhoonidhi implements none of it.

\textbf{Peers implement it.} The Copernicus Data Space publishes 146
server-side processes at a public, unauthenticated endpoint that this project
enumerates live. Bhoonidhi publishes zero. That is not a scoping decision
everyone else also made.

\textbf{The cost is measurable.} Across eight queries the catalogue hands
back 208 scenes where 16 cover the same ground to the same percentage ---
\textbf{92\,\% of everything returned is redundant}, 206\,GB ordered where
16\,GB was required. The user cannot discover this, because discovering it
requires seven manual steps including an NP-hard optimisation and a
reprojection that is easy to get silently wrong. So they order everything.
And because redundancy rises with archive depth while the footprint store
behind this measurement spans only three months, 92\,\% is a floor.

The service built to close that gap implements \textbf{OGC API --- Processes,
Part 1: Core} over the footprint geometry the catalogue already stores, with
six operations spanning both footprint geometry and the raster values beneath
it, synchronous and asynchronous execution, a conformant job API, and all
areas measured in a per-request equal-area projection so that
km\textsuperscript{2} means km\textsuperscript{2}. All 100 compliance tests
pass, and the suite verifies the mathematics --- that covered plus gap
reconstructs the AOI, that an intersection never exceeds its operands, that
greedy coverage is monotonic --- rather than only the protocol shape. Three
real defects were found during development, two of them by those invariant
tests, and each is now guarded.

The single sharpest result is the smallest one. Thirty scenes match a
district-sized area of interest. \textbf{Two cover it.} Finding that out takes
this service nine milliseconds, and takes Bhoonidhi's users an afternoon they
do not spend --- so instead they download the other twenty-eight.

The gap is a \textbf{capability and architecture} gap, not a failing of the
data or of the people who publish it. NRSC makes a remarkable volume of
Indian Earth observation genuinely available. The argument here is narrow
and, I think, fair: a catalogue that knows every footprint's geometry but will
not compute with it leaves its hardest work to every user, one at a time ---
and that last step, from \emph{which scenes touch my area} to \emph{which
scenes I actually need}, is the one that turns a catalogue into
infrastructure.
